
Write a complete proof of the Euclidean property of integers:
If a and b are integers such that a ≥ 0 and b > 0,
then there exists integers q and r such that q ≥ 0 and
a = bq + r, 0 ≤ r < b


We will prove the above statement use the well-ordering principle. For a >= 0, b > 0

a = bq + r, for some q >= 0, and 0 <=  r < b

Consider the set

X = {a - bx | x ∈ N, a - bx > 0}

This set is a nonempty set because x = 0 gives

a - b(0) = a >= 0

Thus a ∈ X

Since X is a nonempty set with non-negative numbers, the WOP says that X has a least element. Let this least element be r. Therefore, for some integer x >= 0

r = a - bx

rearanging gives

a = bx + r

also since r ∈ X

r >= 0

Now, I am going to show that r < b. Assume on the contrary that r >= b, then

r - b >= 0

but r - b = (a - bx) - b = a - b(x + 1)

since x >= 0 we have

x + 1 >= 0

Therefore, r - b ∈ X. However, r - b < r, because b > 0. This is a contradiction because of the fact r was the least element of X. So, our assumption was false and r < b.

Combining everything, we have found integers q, r such that

q >= 0, a = bq + r, 0 <= r < b

Therefore, by the WOP, the Euclidean property of integers holds
